\documentclass[10pt,twocolumn]{IEEEtran}

% ── packages ──────────────────────────────────────────────────────────────────
\usepackage[T1]{fontenc}
\usepackage[utf8]{inputenc}
\usepackage{booktabs}
\usepackage{amsmath}
\usepackage{graphicx}
\usepackage{xcolor}
\usepackage{hyperref}
\usepackage{microtype}
\usepackage{siunitx}
\sisetup{detect-weight=true, mode=text}

% ── convenience ───────────────────────────────────────────────────────────────
\newcommand{\tbd}{\textcolor{gray}{--}}
\newcommand{\na}{\textcolor{gray}{n/a}}
\newcommand{\sect}[1]{\vspace{2pt}\noindent\textbf{#1.}}

% ── graceful fallbacks (overridden by \input below) ───────────────────────────
\providecommand{\whuGAITCnnTestAcc}{\tbd}
\providecommand{\whuGAITAuthTestAcc}{\tbd}
\providecommand{\whuGAITAuthTestAUC}{\tbd}
\providecommand{\whuGAITAuthSep}{\tbd}
\providecommand{\whuGAITAuthOverfitAcc}{\tbd}
\providecommand{\uciharCnnTestAcc}{\tbd}
\providecommand{\uciharAuthTestAcc}{\tbd}
\providecommand{\uciharAuthTestAUC}{\tbd}
\providecommand{\uciharAuthSep}{\tbd}
\providecommand{\uciharAuthOverfitAcc}{\tbd}
\providecommand{\wisdmCnnTestAcc}{\tbd}
\providecommand{\wisdmAuthTestAcc}{\tbd}
\providecommand{\wisdmAuthTestAUC}{\tbd}
\providecommand{\wisdmAuthSep}{\tbd}
\providecommand{\wisdmAuthOverfitAcc}{\tbd}
\providecommand{\combinedCnnTestAcc}{\tbd}
\providecommand{\combinedAuthTestAcc}{\tbd}
\providecommand{\combinedAuthTestAUC}{\tbd}
\providecommand{\combinedAuthSep}{\tbd}
\providecommand{\combinedAuthOverfitAcc}{\tbd}
\providecommand{\whuGAITSimplDAUC}{\tbd}
\providecommand{\whuGAITSimplDDeltaGap}{\tbd}
\providecommand{\uciharSimplDAUC}{\tbd}
\providecommand{\uciharSimplDDeltaGap}{\tbd}
\providecommand{\wisdmSimplDAUC}{\tbd}
\providecommand{\wisdmSimplDDeltaGap}{\tbd}
\providecommand{\combinedSimplDAUC}{\tbd}
\providecommand{\combinedSimplDDeltaGap}{\tbd}
\providecommand{\whuGAITGreyLiraRawAUC}{\tbd}
\providecommand{\whuGAITGreyLiraLfAUC}{\tbd}
\providecommand{\whuGAITGreyLiraMfAUC}{\tbd}
\providecommand{\uciharGreyLiraRawAUC}{\tbd}
\providecommand{\uciharGreyLiraLfAUC}{\tbd}
\providecommand{\uciharGreyLiraMfAUC}{\tbd}
\providecommand{\wisdmGreyLiraRawAUC}{\tbd}
\providecommand{\wisdmGreyLiraLfAUC}{\tbd}
\providecommand{\wisdmGreyLiraMfAUC}{\tbd}
\providecommand{\combinedGreyLiraRawAUC}{\tbd}
\providecommand{\combinedGreyLiraLfAUC}{\tbd}
\providecommand{\combinedGreyLiraMfAUC}{\tbd}
\providecommand{\whuGAITWarmLiraRawAUC}{\tbd}
\providecommand{\whuGAITWarmLiraLfAUC}{\tbd}
\providecommand{\whuGAITWarmLiraMfAUC}{\tbd}
\providecommand{\uciharWarmLiraRawAUC}{\tbd}
\providecommand{\uciharWarmLiraLfAUC}{\tbd}
\providecommand{\uciharWarmLiraMfAUC}{\tbd}
\providecommand{\wisdmWarmLiraRawAUC}{\tbd}
\providecommand{\wisdmWarmLiraLfAUC}{\tbd}
\providecommand{\wisdmWarmLiraMfAUC}{\tbd}
\providecommand{\combinedWarmLiraRawAUC}{\tbd}
\providecommand{\combinedWarmLiraLfAUC}{\tbd}
\providecommand{\combinedWarmLiraMfAUC}{\tbd}
\providecommand{\whuGAITColdLiraRawAUC}{\tbd}
\providecommand{\whuGAITColdLiraLfAUC}{\tbd}
\providecommand{\whuGAITColdLiraMfAUC}{\tbd}
\providecommand{\uciharColdLiraRawAUC}{\tbd}
\providecommand{\uciharColdLiraLfAUC}{\tbd}
\providecommand{\uciharColdLiraMfAUC}{\tbd}
\providecommand{\wisdmColdLiraRawAUC}{\tbd}
\providecommand{\wisdmColdLiraLfAUC}{\tbd}
\providecommand{\wisdmColdLiraMfAUC}{\tbd}
\providecommand{\combinedColdLiraRawAUC}{\tbd}
\providecommand{\combinedColdLiraLfAUC}{\tbd}
\providecommand{\combinedColdLiraMfAUC}{\tbd}
\providecommand{\whuGAITComboASR}{\tbd}
\providecommand{\whuGAITBudgetRatio}{\tbd}
\providecommand{\whuGAITNCombos}{\tbd}
\providecommand{\whuGAITNResistant}{\tbd}
\providecommand{\uciharComboASR}{\tbd}
\providecommand{\uciharBudgetRatio}{\tbd}
\providecommand{\uciharNCombos}{\tbd}
\providecommand{\uciharNResistant}{\tbd}
\providecommand{\wisdmComboASR}{\tbd}
\providecommand{\wisdmBudgetRatio}{\tbd}
\providecommand{\wisdmNCombos}{\tbd}
\providecommand{\wisdmNResistant}{\tbd}
\providecommand{\combinedComboASR}{\tbd}
\providecommand{\combinedBudgetRatio}{\tbd}
\providecommand{\combinedNCombos}{\tbd}
\providecommand{\combinedNResistant}{\tbd}

% ── auto-generated macros ─────────────────────────────────────────────────────
% Auto-generated by compare_models_whuGAIT: DO NOT EDIT MANUALLY
% Last run: 2026-07-28 13:26:26

\providecommand{\whuGAITCnnTestAcc}{}
\renewcommand{\whuGAITCnnTestAcc}{91.44\%}
\providecommand{\whuGAITAuthTestAcc}{}
\renewcommand{\whuGAITAuthTestAcc}{80.12\%}
\providecommand{\whuGAITAuthTrainAcc}{}
\renewcommand{\whuGAITAuthTrainAcc}{91.54\%}
\providecommand{\whuGAITAuthTestAUC}{}
\renewcommand{\whuGAITAuthTestAUC}{0.8698}
\providecommand{\whuGAITAuthOverfitAcc}{}
\renewcommand{\whuGAITAuthOverfitAcc}{+11.42pp}
\providecommand{\whuGAITAuthSameMean}{}
\renewcommand{\whuGAITAuthSameMean}{0.1480}
\providecommand{\whuGAITAuthDiffMean}{}
\renewcommand{\whuGAITAuthDiffMean}{0.6710}
\providecommand{\whuGAITAuthSep}{}
\renewcommand{\whuGAITAuthSep}{0.5220}

% Auto-generated by compare_models_ucihar: DO NOT EDIT MANUALLY
% Last run: 2026-07-28 13:26:26

\providecommand{\uciharCnnTestAcc}{}
\renewcommand{\uciharCnnTestAcc}{85.34\%}
\providecommand{\uciharAuthTestAcc}{}
\renewcommand{\uciharAuthTestAcc}{68.83\%}
\providecommand{\uciharAuthTrainAcc}{}
\renewcommand{\uciharAuthTrainAcc}{81.61\%}
\providecommand{\uciharAuthTestAUC}{}
\renewcommand{\uciharAuthTestAUC}{0.7388}
\providecommand{\uciharAuthOverfitAcc}{}
\renewcommand{\uciharAuthOverfitAcc}{+12.78pp}
\providecommand{\uciharAuthSameMean}{}
\renewcommand{\uciharAuthSameMean}{0.4290}
\providecommand{\uciharAuthDiffMean}{}
\renewcommand{\uciharAuthDiffMean}{0.7200}
\providecommand{\uciharAuthSep}{}
\renewcommand{\uciharAuthSep}{0.2900}

% Auto-generated by compare_models_wisdm: DO NOT EDIT MANUALLY
% Last run: 2026-07-28 13:26:26

\providecommand{\wisdmCnnTestAcc}{}
\renewcommand{\wisdmCnnTestAcc}{96.29\%}
\providecommand{\wisdmAuthTestAcc}{}
\renewcommand{\wisdmAuthTestAcc}{91.87\%}
\providecommand{\wisdmAuthTrainAcc}{}
\renewcommand{\wisdmAuthTrainAcc}{98.02\%}
\providecommand{\wisdmAuthTestAUC}{}
\renewcommand{\wisdmAuthTestAUC}{0.9577}
\providecommand{\wisdmAuthOverfitAcc}{}
\renewcommand{\wisdmAuthOverfitAcc}{+6.15pp}
\providecommand{\wisdmAuthSameMean}{}
\renewcommand{\wisdmAuthSameMean}{0.0870}
\providecommand{\wisdmAuthDiffMean}{}
\renewcommand{\wisdmAuthDiffMean}{0.8790}
\providecommand{\wisdmAuthSep}{}
\renewcommand{\wisdmAuthSep}{0.7920}

% Auto-generated by compare_models_combined: DO NOT EDIT MANUALLY
% Last run: 2026-08-06 21:44:47

\providecommand{\combinedCnnTestAcc}{}
\renewcommand{\combinedCnnTestAcc}{—}
\providecommand{\combinedAuthTestAcc}{}
\renewcommand{\combinedAuthTestAcc}{75.72\%}
\providecommand{\combinedAuthTrainAcc}{}
\renewcommand{\combinedAuthTrainAcc}{82.95\%}
\providecommand{\combinedAuthTestAUC}{}
\renewcommand{\combinedAuthTestAUC}{0.8378}
\providecommand{\combinedAuthOverfitAcc}{}
\renewcommand{\combinedAuthOverfitAcc}{+7.23pp}
\providecommand{\combinedAuthSameMean}{}
\renewcommand{\combinedAuthSameMean}{0.3430}
\providecommand{\combinedAuthDiffMean}{}
\renewcommand{\combinedAuthDiffMean}{0.7370}
\providecommand{\combinedAuthSep}{}
\renewcommand{\combinedAuthSep}{0.3940}

% Auto-generated by compare_mia_whuGAIT: DO NOT EDIT MANUALLY
% Last run: 2026-09-08 12:25:56

\providecommand{\whuGAITSimplDAUC}{}
\renewcommand{\whuGAITSimplDAUC}{0.815}
\providecommand{\whuGAITSimplDDeltaGap}{}
\renewcommand{\whuGAITSimplDDeltaGap}{0.233}
\providecommand{\whuGAITSimplDTPRten}{}
\renewcommand{\whuGAITSimplDTPRten}{0.298}
\providecommand{\whuGAITGreyLiraRawAUC}{}
\renewcommand{\whuGAITGreyLiraRawAUC}{0.262}
\providecommand{\whuGAITGreyLiraRawACC}{}
\renewcommand{\whuGAITGreyLiraRawACC}{0.183}
\providecommand{\whuGAITGreyLiraRawTPRten}{}
\renewcommand{\whuGAITGreyLiraRawTPRten}{0.000}
\providecommand{\whuGAITGreyLiraLfAUC}{}
\renewcommand{\whuGAITGreyLiraLfAUC}{0.163}
\providecommand{\whuGAITGreyLiraLfACC}{}
\renewcommand{\whuGAITGreyLiraLfACC}{0.183}
\providecommand{\whuGAITGreyLiraLfTPRten}{}
\renewcommand{\whuGAITGreyLiraLfTPRten}{0.000}
\providecommand{\whuGAITGreyLiraMfAUC}{}
\renewcommand{\whuGAITGreyLiraMfAUC}{0.251}
\providecommand{\whuGAITGreyLiraMfACC}{}
\renewcommand{\whuGAITGreyLiraMfACC}{0.183}
\providecommand{\whuGAITGreyLiraMfTPRten}{}
\renewcommand{\whuGAITGreyLiraMfTPRten}{0.000}
\providecommand{\whuGAITWarmLiraRawAUC}{}
\renewcommand{\whuGAITWarmLiraRawAUC}{0.346}
\providecommand{\whuGAITWarmLiraRawACC}{}
\renewcommand{\whuGAITWarmLiraRawACC}{0.394}
\providecommand{\whuGAITWarmLiraRawTPRten}{}
\renewcommand{\whuGAITWarmLiraRawTPRten}{0.048}
\providecommand{\whuGAITColdLiraRawAUC}{}
\renewcommand{\whuGAITColdLiraRawAUC}{0.709}
\providecommand{\whuGAITColdLiraRawACC}{}
\renewcommand{\whuGAITColdLiraRawACC}{0.817}
\providecommand{\whuGAITColdLiraRawTPRten}{}
\renewcommand{\whuGAITColdLiraRawTPRten}{0.417}
\providecommand{\whuGAITColdLiraLfAUC}{}
\renewcommand{\whuGAITColdLiraLfAUC}{0.686}
\providecommand{\whuGAITColdLiraLfACC}{}
\renewcommand{\whuGAITColdLiraLfACC}{0.808}
\providecommand{\whuGAITColdLiraLfTPRten}{}
\renewcommand{\whuGAITColdLiraLfTPRten}{0.381}
\providecommand{\whuGAITColdLiraMfAUC}{}
\renewcommand{\whuGAITColdLiraMfAUC}{0.743}
\providecommand{\whuGAITColdLiraMfACC}{}
\renewcommand{\whuGAITColdLiraMfACC}{0.827}
\providecommand{\whuGAITColdLiraMfTPRten}{}
\renewcommand{\whuGAITColdLiraMfTPRten}{0.274}

% Auto-generated by compare_mia_ucihar: DO NOT EDIT MANUALLY
% Last run: 2026-09-08 12:25:56

\providecommand{\uciharSimplDAUC}{}
\renewcommand{\uciharSimplDAUC}{0.794}
\providecommand{\uciharSimplDDeltaGap}{}
\renewcommand{\uciharSimplDDeltaGap}{0.188}
\providecommand{\uciharSimplDTPRten}{}
\renewcommand{\uciharSimplDTPRten}{0.333}
\providecommand{\uciharGreyLiraRawAUC}{}
\renewcommand{\uciharGreyLiraRawAUC}{0.889}
\providecommand{\uciharGreyLiraRawACC}{}
\renewcommand{\uciharGreyLiraRawACC}{0.800}
\providecommand{\uciharGreyLiraRawTPRten}{}
\renewcommand{\uciharGreyLiraRawTPRten}{0.714}
\providecommand{\uciharGreyLiraLfAUC}{}
\renewcommand{\uciharGreyLiraLfAUC}{0.815}
\providecommand{\uciharGreyLiraLfACC}{}
\renewcommand{\uciharGreyLiraLfACC}{0.767}
\providecommand{\uciharGreyLiraLfTPRten}{}
\renewcommand{\uciharGreyLiraLfTPRten}{0.667}
\providecommand{\uciharGreyLiraMfAUC}{}
\renewcommand{\uciharGreyLiraMfAUC}{0.862}
\providecommand{\uciharGreyLiraMfACC}{}
\renewcommand{\uciharGreyLiraMfACC}{0.700}
\providecommand{\uciharGreyLiraMfTPRten}{}
\renewcommand{\uciharGreyLiraMfTPRten}{0.571}
\providecommand{\uciharWarmLiraRawAUC}{}
\renewcommand{\uciharWarmLiraRawAUC}{0.661}
\providecommand{\uciharWarmLiraRawACC}{}
\renewcommand{\uciharWarmLiraRawACC}{0.667}
\providecommand{\uciharWarmLiraRawTPRten}{}
\renewcommand{\uciharWarmLiraRawTPRten}{0.000}
\providecommand{\uciharColdLiraRawAUC}{}
\renewcommand{\uciharColdLiraRawAUC}{0.968}
\providecommand{\uciharColdLiraRawACC}{}
\renewcommand{\uciharColdLiraRawACC}{0.700}
\providecommand{\uciharColdLiraRawTPRten}{}
\renewcommand{\uciharColdLiraRawTPRten}{0.905}
\providecommand{\uciharColdLiraLfAUC}{}
\renewcommand{\uciharColdLiraLfAUC}{0.921}
\providecommand{\uciharColdLiraLfACC}{}
\renewcommand{\uciharColdLiraLfACC}{0.700}
\providecommand{\uciharColdLiraLfTPRten}{}
\renewcommand{\uciharColdLiraLfTPRten}{0.667}
\providecommand{\uciharColdLiraMfAUC}{}
\renewcommand{\uciharColdLiraMfAUC}{0.926}
\providecommand{\uciharColdLiraMfACC}{}
\renewcommand{\uciharColdLiraMfACC}{0.700}
\providecommand{\uciharColdLiraMfTPRten}{}
\renewcommand{\uciharColdLiraMfTPRten}{0.762}

% Auto-generated by compare_mia_wisdm: DO NOT EDIT MANUALLY
% Last run: 2026-09-08 12:25:56

\providecommand{\wisdmSimplDAUC}{}
\renewcommand{\wisdmSimplDAUC}{0.952}
\providecommand{\wisdmSimplDDeltaGap}{}
\renewcommand{\wisdmSimplDDeltaGap}{0.136}
\providecommand{\wisdmSimplDTPRten}{}
\renewcommand{\wisdmSimplDTPRten}{0.950}
\providecommand{\wisdmGreyLiraRawAUC}{}
\renewcommand{\wisdmGreyLiraRawAUC}{0.827}
\providecommand{\wisdmGreyLiraRawACC}{}
\renewcommand{\wisdmGreyLiraRawACC}{0.706}
\providecommand{\wisdmGreyLiraRawTPRten}{}
\renewcommand{\wisdmGreyLiraRawTPRten}{0.625}
\providecommand{\wisdmGreyLiraLfAUC}{}
\renewcommand{\wisdmGreyLiraLfAUC}{0.723}
\providecommand{\wisdmGreyLiraLfACC}{}
\renewcommand{\wisdmGreyLiraLfACC}{0.608}
\providecommand{\wisdmGreyLiraLfTPRten}{}
\renewcommand{\wisdmGreyLiraLfTPRten}{0.525}
\providecommand{\wisdmGreyLiraMfAUC}{}
\renewcommand{\wisdmGreyLiraMfAUC}{0.857}
\providecommand{\wisdmGreyLiraMfACC}{}
\renewcommand{\wisdmGreyLiraMfACC}{0.765}
\providecommand{\wisdmGreyLiraMfTPRten}{}
\renewcommand{\wisdmGreyLiraMfTPRten}{0.725}
\providecommand{\wisdmWarmLiraRawAUC}{}
\renewcommand{\wisdmWarmLiraRawAUC}{0.659}
\providecommand{\wisdmWarmLiraRawACC}{}
\renewcommand{\wisdmWarmLiraRawACC}{0.784}
\providecommand{\wisdmWarmLiraRawTPRten}{}
\renewcommand{\wisdmWarmLiraRawTPRten}{0.250}
\providecommand{\wisdmColdLiraRawAUC}{}
\renewcommand{\wisdmColdLiraRawAUC}{0.916}
\providecommand{\wisdmColdLiraRawACC}{}
\renewcommand{\wisdmColdLiraRawACC}{0.784}
\providecommand{\wisdmColdLiraRawTPRten}{}
\renewcommand{\wisdmColdLiraRawTPRten}{0.800}
\providecommand{\wisdmColdLiraLfAUC}{}
\renewcommand{\wisdmColdLiraLfAUC}{0.968}
\providecommand{\wisdmColdLiraLfACC}{}
\renewcommand{\wisdmColdLiraLfACC}{0.784}
\providecommand{\wisdmColdLiraLfTPRten}{}
\renewcommand{\wisdmColdLiraLfTPRten}{0.925}
\providecommand{\wisdmColdLiraMfAUC}{}
\renewcommand{\wisdmColdLiraMfAUC}{0.968}
\providecommand{\wisdmColdLiraMfACC}{}
\renewcommand{\wisdmColdLiraMfACC}{0.804}
\providecommand{\wisdmColdLiraMfTPRten}{}
\renewcommand{\wisdmColdLiraMfTPRten}{0.950}

% Auto-generated by compare_mia_combined: DO NOT EDIT MANUALLY
% Last run: 2026-07-28 13:55:15

\providecommand{\combinedSimplDAUC}{}
\renewcommand{\combinedSimplDAUC}{0.753}
\providecommand{\combinedSimplDDeltaGap}{}
\renewcommand{\combinedSimplDDeltaGap}{0.262}
\providecommand{\combinedGreyLiraRawACC}{}
\renewcommand{\combinedGreyLiraRawACC}{—}
\providecommand{\combinedGreyLiraRawAUC}{}
\renewcommand{\combinedGreyLiraRawAUC}{—}
\providecommand{\combinedGreyLiraLfACC}{}
\renewcommand{\combinedGreyLiraLfACC}{—}
\providecommand{\combinedGreyLiraLfAUC}{}
\renewcommand{\combinedGreyLiraLfAUC}{—}
\providecommand{\combinedGreyLiraMfACC}{}
\renewcommand{\combinedGreyLiraMfACC}{—}
\providecommand{\combinedGreyLiraMfAUC}{}
\renewcommand{\combinedGreyLiraMfAUC}{—}
\providecommand{\combinedWarmLiraRawACC}{}
\renewcommand{\combinedWarmLiraRawACC}{—}
\providecommand{\combinedWarmLiraRawAUC}{}
\renewcommand{\combinedWarmLiraRawAUC}{—}
\providecommand{\combinedWarmLiraLfACC}{}
\renewcommand{\combinedWarmLiraLfACC}{—}
\providecommand{\combinedWarmLiraLfAUC}{}
\renewcommand{\combinedWarmLiraLfAUC}{—}
\providecommand{\combinedWarmLiraMfACC}{}
\renewcommand{\combinedWarmLiraMfACC}{—}
\providecommand{\combinedWarmLiraMfAUC}{}
\renewcommand{\combinedWarmLiraMfAUC}{—}
\providecommand{\combinedColdLiraRawACC}{}
\renewcommand{\combinedColdLiraRawACC}{—}
\providecommand{\combinedColdLiraRawAUC}{}
\renewcommand{\combinedColdLiraRawAUC}{—}
\providecommand{\combinedColdLiraLfACC}{}
\renewcommand{\combinedColdLiraLfACC}{—}
\providecommand{\combinedColdLiraLfAUC}{}
\renewcommand{\combinedColdLiraLfAUC}{—}
\providecommand{\combinedColdLiraMfACC}{}
\renewcommand{\combinedColdLiraMfACC}{—}
\providecommand{\combinedColdLiraMfAUC}{}
\renewcommand{\combinedColdLiraMfAUC}{—}

% Auto-generated by compare_evasion_whuGAIT: DO NOT EDIT MANUALLY
% Last run: 2026-08-06 21:45:02

\providecommand{\whuGAITEpsTarget}{}
\renewcommand{\whuGAITEpsTarget}{24.54}
\providecommand{\whuGAITComboASR}{}
\renewcommand{\whuGAITComboASR}{1.000}
\providecommand{\whuGAITBudgetRatio}{}
\renewcommand{\whuGAITBudgetRatio}{0.056}
\providecommand{\whuGAITNResistant}{}
\renewcommand{\whuGAITNResistant}{0}
\providecommand{\whuGAITNPairs}{}
\renewcommand{\whuGAITNPairs}{3800}
\providecommand{\whuGAITNFullCombos}{}
\renewcommand{\whuGAITNFullCombos}{220}
\providecommand{\whuGAITNCombos}{}
\renewcommand{\whuGAITNCombos}{220}
\providecommand{\whuGAITTrainComboASR}{}
\renewcommand{\whuGAITTrainComboASR}{0.914}
\providecommand{\whuGAITTrainBudgetRatio}{}
\renewcommand{\whuGAITTrainBudgetRatio}{0.269}

% Auto-generated by compare_evasion_ucihar: DO NOT EDIT MANUALLY
% Last run: 2026-07-28 13:26:39

\providecommand{\uciharEpsTarget}{}
\renewcommand{\uciharEpsTarget}{37.47}
\providecommand{\uciharComboASR}{}
\renewcommand{\uciharComboASR}{0.980}
\providecommand{\uciharBudgetRatio}{}
\renewcommand{\uciharBudgetRatio}{0.135}
\providecommand{\uciharNResistant}{}
\renewcommand{\uciharNResistant}{106}
\providecommand{\uciharNPairs}{}
\renewcommand{\uciharNPairs}{5000}
\providecommand{\uciharNFullCombos}{}
\renewcommand{\uciharNFullCombos}{42}
\providecommand{\uciharNCombos}{}
\renewcommand{\uciharNCombos}{72}
\providecommand{\uciharTrainComboASR}{}
\renewcommand{\uciharTrainComboASR}{0.970}
\providecommand{\uciharTrainBudgetRatio}{}
\renewcommand{\uciharTrainBudgetRatio}{0.148}

% Auto-generated by compare_evasion_wisdm: DO NOT EDIT MANUALLY
% Last run: 2026-07-28 13:26:39

\providecommand{\wisdmEpsTarget}{}
\renewcommand{\wisdmEpsTarget}{31.00}
\providecommand{\wisdmComboASR}{}
\renewcommand{\wisdmComboASR}{0.978}
\providecommand{\wisdmBudgetRatio}{}
\renewcommand{\wisdmBudgetRatio}{0.300}
\providecommand{\wisdmNResistant}{}
\renewcommand{\wisdmNResistant}{106}
\providecommand{\wisdmNPairs}{}
\renewcommand{\wisdmNPairs}{5000}
\providecommand{\wisdmNFullCombos}{}
\renewcommand{\wisdmNFullCombos}{88}
\providecommand{\wisdmNCombos}{}
\renewcommand{\wisdmNCombos}{110}
\providecommand{\wisdmTrainComboASR}{}
\renewcommand{\wisdmTrainComboASR}{0.932}
\providecommand{\wisdmTrainBudgetRatio}{}
\renewcommand{\wisdmTrainBudgetRatio}{0.415}

% Auto-generated by compare_evasion_combined: DO NOT EDIT MANUALLY
% Last run: 2026-07-28 13:26:39

\providecommand{\combinedEpsTarget}{}
\renewcommand{\combinedEpsTarget}{24.11}
\providecommand{\combinedComboASR}{}
\renewcommand{\combinedComboASR}{0.962}
\providecommand{\combinedBudgetRatio}{}
\renewcommand{\combinedBudgetRatio}{0.116}
\providecommand{\combinedNResistant}{}
\renewcommand{\combinedNResistant}{125}
\providecommand{\combinedNPairs}{}
\renewcommand{\combinedNPairs}{4000}
\providecommand{\combinedNFullCombos}{}
\renewcommand{\combinedNFullCombos}{146}
\providecommand{\combinedNCombos}{}
\renewcommand{\combinedNCombos}{182}
\providecommand{\combinedTrainComboASR}{}
\renewcommand{\combinedTrainComboASR}{0.944}
\providecommand{\combinedTrainBudgetRatio}{}
\renewcommand{\combinedTrainBudgetRatio}{0.454}


% ── metadata ──────────────────────────────────────────────────────────────────
\title{\large\textbf{Gait MIA Pipeline -- Results Summary}\\[2pt]
       \normalsize Laboratory Note \quad v0.3}

\author{Matteo Campagnaro\\
        \small Supervisor: Prof.\ Simone Milani, University of Padova\\
        \small \today}

\begin{document}
\maketitle
\thispagestyle{plain}

\begin{abstract}
This note collects numerical results from the gait-based membership inference
(MIA) pipeline across four configurations: within-dataset runs on whuGAIT,
UCI-HAR, and WISDM, and a cross-dataset run (Combined) that reuses the whuGAIT
CNN encoder with an LSTM authenticator retrained on subjects from all three
sources. Values are reported without selective filtering. Observations are
tentative and subject to revision as experimental conditions change.
\end{abstract}

% ══════════════════════════════════════════════════════════════════════════════
\section{Experimental Setup}
% ══════════════════════════════════════════════════════════════════════════════

\begin{table}[h]
\centering
\caption{Subject counts and authentication pair volumes (quick-mode).}
\label{tab:setup}
\setlength{\tabcolsep}{3.5pt}
\begin{tabular}{lrrrr}
\toprule
Config   & Members & Non-mbr & Train pairs & Test pairs \\
\midrule
whuGAIT  & 98  & 20 & 66\,542 & 7\,600  \\
UCI-HAR  & 21  &  9 & 60\,000 & 10\,000 \\
WISDM    & 40  & 11 & 60\,000 & 10\,000 \\
Combined & 159 & 20 & 72\,000 &  8\,000 \\
\bottomrule
\end{tabular}
\end{table}

\sect{whuGAIT} Subjects 21--118 are members; 1--20 are non-members. Pairs come
from Dataset~\#5 (pre-segmented, step-synchronised windows). A correlation
asymmetry in the pair structure confounds LiRA (Section~\ref{sec:mia}).

\sect{Combined} whuGAIT subjects 21--118 are members; UCI-HAR train-split (21)
and WISDM 1--40 contribute cross-dataset members; UCI-HAR test-split (9) and
WISDM 41--51 are non-members. The CNN encoder is frozen from the whuGAIT
checkpoint; only the LSTM is retrained. Pairs are built within each source
separately before concatenation.

% ══════════════════════════════════════════════════════════════════════════════
\section{CNN Encoder}
\label{sec:cnn}
% ══════════════════════════════════════════════════════════════════════════════

\begin{table}[h]
\centering
\caption{CNN identification accuracy (top-1, test split, 2 epochs).}
\label{tab:cnn}
\setlength{\tabcolsep}{4pt}
\begin{tabular}{lrl}
\toprule
Config   & Test acc           & Notes \\
\midrule
whuGAIT  & \whuGAITCnnTestAcc & 98-class \\
UCI-HAR  & \uciharCnnTestAcc  & 6-class  \\
WISDM    & \wisdmCnnTestAcc   & 51-class \\
Combined & \na                & frozen encoder \\
\bottomrule
\end{tabular}
\end{table}

% ══════════════════════════════════════════════════════════════════════════════
\section{Authentication Model}
\label{sec:auth}
% ══════════════════════════════════════════════════════════════════════════════

\begin{table}[h]
\centering
\caption{CNN+LSTM authenticator on held-out test pairs. Sep $=
\bar{P}_{\text{imp}} - \bar{P}_{\text{gen}}$. Overfit $=$ train acc $-$ test
acc at the best checkpoint epoch.}
\label{tab:auth}
\setlength{\tabcolsep}{3pt}
\begin{tabular}{lrrrr}
\toprule
Config   & Test acc              & AUC                  & Sep                  & Overfit \\
\midrule
whuGAIT  & \whuGAITAuthTestAcc   & \whuGAITAuthTestAUC  & \whuGAITAuthSep      & \whuGAITAuthOverfitAcc  \\
UCI-HAR  & \uciharAuthTestAcc    & \uciharAuthTestAUC   & \uciharAuthSep       & \uciharAuthOverfitAcc   \\
WISDM    & \wisdmAuthTestAcc     & \wisdmAuthTestAUC    & \wisdmAuthSep        & \wisdmAuthOverfitAcc    \\
Combined & \combinedAuthTestAcc  & \combinedAuthTestAUC & \combinedAuthSep     & \combinedAuthOverfitAcc \\
\bottomrule
\end{tabular}
\end{table}

The Overfit column reflects the train-test accuracy gap at the best saved epoch.
The Combined run shows the largest gap (\combinedAuthOverfitAcc), which may
indicate the LSTM is memorising subject-level patterns from training.
Fig.~\ref{fig:overfitting} shows per-epoch accuracy and loss curves for each
configuration; the gap between train and test tracks overfitting directly.

\begin{figure*}[t]
\centering
\includegraphics[width=\textwidth]{../../results/analysis/models_overfitting_curves.png}
\caption{Authentication model (CNN+LSTM) training curves per epoch. Top row:
         accuracy; bottom row: cross-entropy loss. Dashed = test, solid = train.
         Vertical dotted line marks the best-checkpoint epoch. Quick-mode uses
         2 epochs so curves have not necessarily converged.}
\label{fig:overfitting}
\end{figure*}

% ══════════════════════════════════════════════════════════════════════════════
\section{Membership Inference}
\label{sec:mia}
% ══════════════════════════════════════════════════════════════════════════════

The per-subject membership signal is the delta
\begin{equation}
  \delta(s) = \bar{P}_{\text{imp}}(s) - \bar{P}_{\text{gen}}(s),
  \label{eq:delta}
\end{equation}
where $\bar{P}_{\text{imp}}$ and $\bar{P}_{\text{gen}}$ are mean
$P(\text{different person})$ over impostor and genuine pairs for subject~$s$.
Fig.~\ref{fig:deltas} shows the per-subject distributions; a larger horizontal
separation indicates a stronger signal.

\begin{figure*}[t]
\centering
\includegraphics[width=\textwidth]{../../results/analysis/mia_nb04_deltas.png}
\caption{Per-subject delta distributions for members (blue) and non-members
         (orange) across all four configurations.}
\label{fig:deltas}
\end{figure*}

\begin{table}[h]
\centering
\caption{Simple-delta MIA. $\Delta$ gap $= \bar\delta_M - \bar\delta_{NM}$.}
\label{tab:mia-signal}
\setlength{\tabcolsep}{4pt}
\begin{tabular}{lrr}
\toprule
Config   & $\Delta$ gap            & AUC \\
\midrule
whuGAIT  & \whuGAITSimplDDeltaGap  & \whuGAITSimplDAUC  \\
UCI-HAR  & \uciharSimplDDeltaGap   & \uciharSimplDAUC   \\
WISDM    & \wisdmSimplDDeltaGap    & \wisdmSimplDAUC    \\
Combined & \combinedSimplDDeltaGap & \combinedSimplDAUC \\
\bottomrule
\end{tabular}
\end{table}

\subsection{LiRA threat model comparison}

Table~\ref{tab:lira-threat} compares raw-delta AUC across three threat models.
Grey-box initialises shadows from the actual target weights; warm uses an
independently trained surrogate; cold starts from random Xavier init. All use
member-aware splits ($K{=}4$ quick-mode shadows).

\begin{table}[h]
\centering
\caption{LiRA AUC (raw-delta) by threat model. ``--'' indicates no valid
         estimates ($K{=}4 < 5$ required per subject).}
\label{tab:lira-threat}
\setlength{\tabcolsep}{3pt}
\begin{tabular}{lrrrr}
\toprule
Config   & Simple-D            & Grey                    & Warm                    & Cold \\
\midrule
whuGAIT  & \whuGAITSimplDAUC   & \whuGAITGreyLiraRawAUC  & \whuGAITWarmLiraRawAUC  & \whuGAITColdLiraRawAUC  \\
UCI-HAR  & \uciharSimplDAUC    & \uciharGreyLiraRawAUC   & \uciharWarmLiraRawAUC   & \uciharColdLiraRawAUC   \\
WISDM    & \wisdmSimplDAUC     & \wisdmGreyLiraRawAUC    & \wisdmWarmLiraRawAUC    & \wisdmColdLiraRawAUC    \\
Combined & \combinedSimplDAUC  & \combinedGreyLiraRawAUC & \combinedWarmLiraRawAUC & \combinedColdLiraRawAUC \\
\bottomrule
\end{tabular}
\end{table}

\subsection{LiRA delta variant comparison}

Three delta formulations are evaluated for the cold-start threat model, which
provides the cleanest comparison as it requires no knowledge of the target:
\begin{align*}
\delta_\text{raw} &= \bar{p} - \bar{q}, \\
\delta_\text{lf}  &= \overline{\mathrm{logit}(p)} - \overline{\mathrm{logit}(q)}, \\
\delta_\text{mf}  &= \mathrm{logit}(\bar{p}) - \mathrm{logit}(\bar{q}).
\end{align*}

\begin{table}[h]
\centering
\caption{Cold-start LiRA AUC by delta variant. Simple-D shown for reference.}
\label{tab:lira-variants}
\setlength{\tabcolsep}{3pt}
\begin{tabular}{lrrrr}
\toprule
Config   & Simple-D            & Cold-raw                & Cold-lf                & Cold-mf \\
\midrule
whuGAIT  & \whuGAITSimplDAUC   & \whuGAITColdLiraRawAUC  & \whuGAITColdLiraLfAUC  & \whuGAITColdLiraMfAUC  \\
UCI-HAR  & \uciharSimplDAUC    & \uciharColdLiraRawAUC   & \uciharColdLiraLfAUC   & \uciharColdLiraMfAUC   \\
WISDM    & \wisdmSimplDAUC     & \wisdmColdLiraRawAUC    & \wisdmColdLiraLfAUC    & \wisdmColdLiraMfAUC    \\
Combined & \combinedSimplDAUC  & \combinedColdLiraRawAUC & \combinedColdLiraLfAUC & \combinedColdLiraMfAUC \\
\bottomrule
\end{tabular}
\end{table}

Fig.~\ref{fig:roc} shows ROC curves comparing Simple-delta, grey-box, and
cold-start LiRA across all configurations. The per-dataset subplots make the
threat-model comparison and the whuGAIT inversion effect directly visible.

\begin{figure*}[t]
\centering
\includegraphics[width=\textwidth]{../../results/analysis/mia_roc_comparison.png}
\caption{MIA ROC curves per configuration (quick-mode, $K{=}4$ shadows).
         Solid: Simple-delta. Dashed: grey-box LiRA (raw, lf).
         Dotted: cold-start LiRA (raw, lf). Combined shows Simple-delta only
         because $K{=}4$ shadows produce no valid LiRA estimates. The whuGAIT
         grey-box curves fall below the diagonal due to the D5 pair-correlation
         confound.}
\label{fig:roc}
\end{figure*}

% ══════════════════════════════════════════════════════════════════════════════
\section{Evasion Attack}
\label{sec:evasion}
% ══════════════════════════════════════════════════════════════════════════════

\begin{table}[h]
\centering
\caption{PGD impersonation attack on test impostor pairs. Budget ratio:
         $\bar\varepsilon_{\min}/\varepsilon_{\text{target}}$ where
         $\varepsilon_{\text{target}}$ is the median intra-subject $\ell_2$
         distance. Resistant: combos unsolved across the full $\varepsilon$ grid.}
\label{tab:evasion}
\setlength{\tabcolsep}{3.5pt}
\begin{tabular}{lrrrr}
\toprule
Config   & Combos              & ASR                 & Budget ratio         & Resistant \\
\midrule
whuGAIT  & \whuGAITNCombos     & \whuGAITComboASR    & \whuGAITBudgetRatio  & \whuGAITNResistant  \\
UCI-HAR  & \uciharNCombos      & \uciharComboASR     & \uciharBudgetRatio   & \uciharNResistant   \\
WISDM    & \wisdmNCombos       & \wisdmComboASR      & \wisdmBudgetRatio    & \wisdmNResistant    \\
Combined & \combinedNCombos    & \combinedComboASR   & \combinedBudgetRatio & \combinedNResistant \\
\bottomrule
\end{tabular}
\end{table}

Budget ratios below 1.0 indicate the required perturbation falls within natural
intra-subject gait variation. The whuGAIT ratio (\whuGAITBudgetRatio) is the
lowest, while Combined has the highest number of resistant combos
(\combinedNResistant vs 0 for whuGAIT), which may reflect added decision-boundary
complexity from cross-dataset training.

Fig.~\ref{fig:evasion-summary} shows ASR and budget ratio per dataset;
Fig.~\ref{fig:evasion-eps} shows the distribution of minimum perturbation budgets
($\varepsilon_{\min}$) across successful pairs; Fig.~\ref{fig:evasion-traintest}
compares test-split and training-split attack outcomes.

For whuGAIT, spectral analysis (NB06c) shows perturbation energy concentrated
in the gait frequency band (0--5\,Hz) and a Mann-Whitney test on the
high-frequency ratio gives $p{=}0.82$, indicating no detectable high-frequency
artefact under this detector.

\begin{figure*}[t]
\centering
\includegraphics[width=\textwidth]{../../results/analysis/evasion_summary_bars.png}
\caption{Evasion attack summary. Left: budget ratio
         ($\bar\varepsilon_{\min}/\varepsilon_{\text{target}}$) per dataset;
         values below 1.0 indicate the attack stays within natural gait variation.
         Right: attack success rate (ASR) at the combo level.}
\label{fig:evasion-summary}
\end{figure*}

\begin{figure*}[t]
\centering
\includegraphics[width=\textwidth]{../../results/analysis/evasion_eps_distribution.png}
\caption{Distribution of minimum perturbation budgets $\varepsilon_{\min}$
         across successfully attacked pairs. Narrower distributions indicate
         more consistent attack difficulty across the test set.}
\label{fig:evasion-eps}
\end{figure*}

\begin{figure*}[t]
\centering
\includegraphics[width=\textwidth]{../../results/analysis/evasion_train_vs_test.png}
\caption{Train-split vs test-split evasion results. Training pairs are generally
         harder to attack (fewer epochs, different pair structure), so the
         train-split ASR and budget ratio tend to differ from the test-split.}
\label{fig:evasion-traintest}
\end{figure*}

% ══════════════════════════════════════════════════════════════════════════════
\section{Summary Observations}
% ══════════════════════════════════════════════════════════════════════════════

\begin{itemize}\setlength\itemsep{1pt}
  \item Simple-delta AUC exceeds 0.5 in all four configurations, indicating
        some degree of memorisation by the authenticator in each case.
  \item The D5 pair-correlation confound causes grey-box and warm-start LiRA
        to invert for whuGAIT (AUC $< 0.5$); cold-start is less affected.
        The effect is visible in the ROC curves (Fig.~\ref{fig:roc}).
  \item Cold-start LiRA tends to outperform warm-start across datasets, a
        pattern that may indicate overfitting in the surrogate training phase.
  \item Among the delta variants, raw and mean-first perform similarly for
        cold-start; logit-first is generally weaker.
  \item The Combined run shows the largest authentication overfit gap
        (\combinedAuthOverfitAcc) and moderate simple-delta AUC
        (\combinedSimplDAUC), consistent with proportional underrepresentation
        of UCI-HAR pairs ($\approx$8\% of the training set).
  \item Evasion budget ratios below 0.3 across all configurations suggest the
        PGD perturbation tends to stay within natural intra-subject variation,
        though quick-mode ($K_\text{PGD}{=}5$) results may differ from
        full-mode behaviour.
\end{itemize}

% ══════════════════════════════════════════════════════════════════════════════
\section*{Conditions}
% ══════════════════════════════════════════════════════════════════════════════

\noindent
\textbf{Hardware:} Apple M-series (MPS)\quad
\textbf{Framework:} PyTorch 2.x\quad
\textbf{Seed:} 42 (splits), 0 (shadow init)\quad
\textbf{Mode:} quick (CNN 2 ep, LSTM 2 ep, $K_\text{LiRA}{=}4$, $K_\text{PGD}{=}5$)

\vspace{4pt}
\noindent\textit{LiRA results for Combined are unavailable at $K{=}4$ shadows.
All results are expected to change under full-mode training.}

\end{document}
